\documentclass[11pt]{article}
\usepackage[utf8]{inputenc}
\usepackage[spanish]{babel}
\usepackage[margin=2.5cm]{geometry}
\usepackage{enumitem}
\usepackage{titlesec}

\title{1º Idea: Misterio/Suspenso — Primera Persona}
\author{}
\date{}

\begin{document}
\maketitle

\begin{enumerate}[label=\arabic*)]
    \item \textbf{Lente cromático:} sostener una linterna/vela/candelabro con 3 colores (Rojo, Azul, violeta). Cada carga hace visible distintas cosas de la habitación: cables ocultos, mensajes o mecanismos invisibles.

    \item Innovar en la exploración, fijándote no dejarte nada de un color.

    \item Abrir una caja fuerte oculta en la pieza del final de una mansión para recuperar un tesoro y salir de la mansión.

    \item Juego escape y misterio en primera persona en una sola habitación detallada (Escape Room / Mystery Exploration).

    \item El jugador mira un cuadro aparentemente normal; selecciona un color, aparecen números fosforescentes sobre el marco y huellas luminosas en el suelo que conducen a un libro falso en la estantería.

    \item Las pistas y cerraduras no se ven con luz normal; el avance depende solamente de contrastar lo que revela cada filtro o color.
\end{enumerate}

\section*{Diseño de niveles}
Mansión con un hall, escaleras grandes para subir a la habitación del medio bloqueada. Dos pisos: dos piezas por piso: una a la derecha y una a la izquierda.

\section*{Estética}
Mansión abandonada. Indicar por qué usar tal luz en cual momento dado. ¿Mezclar luces?

\end{document}
